\externaldocument[nocite]{build/text/chapters/our-games}[url=text/chapters/our-games]
%===============================================================
\chapter{Launching and Testing the Programme}\label{chap:launch-testing}
%===============================================================

\begin{chapterabstract}
	To ensure the long-term viability of the student game incubation programme, it needs to be launched with broad outreach to effectively validate the chosen approach with direct user feedback and adaptation on the go. This chapter describes the grafit.games website—launched to inform, register, and engage potential participants—and the campaign, utilising posters put up at FIT and in the relevant Discord channels to raise awareness. User sentiment testing has begun, aiming to understand developer priorities and expectations. For pilot testing, participants will be recruited through direct outreach to experienced game authors, promotions at events (e.g., GameJams), and open applications at CTU. The year-long pilot phase will be followed by a structured evaluation to assess participant benefit, content relevance, delivery methods, and operational feasibility.
\end{chapterabstract}

To ensure long-term viability of the incubation programme, I and the future staff of Grafit.games aim to launch it with sizable outreach. I want to gather broad up-front feedback in the community and to start widely inviting participants after launch. This approach will help verify the assumptions I made when designing the mechanisms and make iterative improvements.

%===============================================================
\section{User Sentiment Testing}
Understanding what student developers need, prioritise, and how they perceive the programme is essential to its successful launch. Therefore, I have started gathering opinions and conducting user sentiment testing.

%---------------------------------------------------------------
\subsection{Presenting through a Website}
To direct outreach, I have created a dedicated website designed to fit well into the Grafit design language and described in the appendix \ref{ap:web}.
\begin{figure}[H]
    \includegraphics[width=0.9\textwidth]{web_landing.png}
    \caption{The grafit.games landing page}
    \label{fig:web_landing}
\end{figure}

This platform serves several key purposes:
\begin{itemize}
    \item It informs potential participants about the goals, structure, and benefits of the programme:
    \begin{figure}[H]
        \includegraphics[width=0.9\textwidth]{web_our-proposition.png}
        \caption{Screenshot from grafit.games - Our Proposition}
        \label{fig:web_our-proposition}
    \end{figure}
    \item It provides updates on the programme’s progress and outlines upcoming events or changes:
    \begin{figure}[H]
        \includegraphics[width=0.9\textwidth]{web_current-events.png}
        \caption{Screenshot from grafit.games - Current}
        \label{fig:web_current-events}
    \end{figure}
    \item It allows interested individuals to register or express their intent to contribute. I am collecting a name, email (required), Discord username, game description, input on what the interviewee would like to see from Grafit.games, whether they would like to participate in the preparation of the programme and whether they would like to be informed about updates:
    \begin{figure}[H]    
        \includegraphics[width=0.9\textwidth]{web_get-notified.png}
        \caption{Screenshot from grafit.games - Get Notified}
        \label{fig:web_get-notified}
    \end{figure}
\end{itemize}

The website serves as a communication channel, evokes credibility, and promises progress to potential participants and backers.
        
%---------------------------------------------------------------
\subsection{Campaign}
\begin{figure}[H]
    To direct interested individuals to the website, I’ve put up posters on the premises of our faculty and in the relevant Discord servers. The performance of the campaign is tracked using bitly\foornote{Accesible through bitly.com}.
    
    \vspace{1em}
    \includegraphics[width=0.8\textwidth]{posterPrint.png}
    \caption{The grafit.games poster}
    \label{fig:poster}
\end{figure}

%===============================================================
\section{Pilot Run}
For a pilot run starting in the winter semester of the 2025/2026 academic year, I will recruit a group of developers representative of the target audience, aiming to gain a broad range of feedback. I will focus on multiple sources for enrolments:
\begin{itemize}
    \item \textbf{Direct Outreach to Game Authors:} I will directly approach creators of exemplary games---some of which have been referenced in section \ref{sec:example-games}. These individuals have firsthand experience with the development process within the academic environment, can provide informed feedback and, crucially, possess a game in the late stages of development that could be commercialised.
    \item \textbf{Engagement at Events:} I will invite competitors at hackathons such as GameJam and advertise the activities of the incubator at promotional events.
    \item \textbf{Open Applications at CTU:} Whilst I mostly expect interest from students at FIT, the programme will be available to all associates of CTU---students, staff, and alumni alike.
\end{itemize}
I will continue utilising the campaign for sentiment testing and gather direct feedback from members using forms.

%===============================================================
\section{Evaluating the Support Mechanism}
I want Grafit.games to improve as much as possible throughout a year-long pilot run. For the consecutive summer, however, a structured evaluation phase is planned. I will assess the strengths, weaknesses, impact, and sustainability of the programme and make modifications if necessary.

%---------------------------------------------------------------
\subsection{Evaluation Criteria}
These criteria will be included in the participant questionnaires and then considered in the evaluation:\begin{itemize}
    \item \textbf{Participant Benefit:} I aim to determine whether participants found the programme helpful and relevant to their needs.
    \item \textbf{Content Evaluation:} I will assess whether the teaching materials cover necessary technical, legal, and business topics without overwhelming participants or distracting from their studies.
    \item \textbf{Delivery Optimization:} I will evaluate whether the chosen mediums (e.g., workshops, online tutorials, consultations) are effective in conveying the content.
\end{itemize}

Additionally, I will determine and consider:
\begin{itemize}
    \item \textbf{Staffing Feasibility:} I will review whether the project will be capable of attaining sufficient staffing to operate the incubator.
    \item \textbf{Financial Feasibility:} I will review the expenses needed to run the incubator and whether it can attract financial backers.
    \item \textbf{Feasibility of the Selected Legal Form:} I will evaluate whether the legal and operational framework poses any risk to the programme’s longevity.
\end{itemize}

%---------------------------------------------------------------
\subsection{Procedure After Evaluation}
Based on the findings from the year-long operations, I and the staff of Grafit.games will be prepared to take decisive action:
\begin{itemize}
    \item \textbf{Adjustments in the Teaching:} In the case that 30\% of the respondents considered the programme unhelpful, or if major insufficiencies in the taught information---whether it is the contents, structure, or delivery---have been identified by evaluating the questionnaires, then the programme will be restructured for a consecutive run the next year.
    \item \textbf{Changing the Legal Structure:} If it is determined that the legal, executive, or operational framework does not fit the programme’s needs, it will be replaced.
    \item \textbf{Termination and Redesign:} Should the programme be deemed unhelpful by a majority (50\%) of the respondents, or should it be deemed unsustainable in its current form, I and the staff will cancel it and---hopefully---propose an alternative solution that better fits the needs of the community.
\end{itemize}

I am confident that this flexibility will enable the initiative to reach the quality and impact I envision.
